% CSE 719 — Distributed & Cloud Computing
% Lecture 1: Distributed System Fundamentals — Model Answers (2020–2021)
% University of Chittagong, 8th Semester B.Sc. Engineering
% Note: 2022 and 2024 papers contain no Lecture-1 questions.
% Compile: pdflatex Fundamentals_Solutions.tex (twice for TOC)

\documentclass[12pt, a4paper]{article}

\usepackage[left=2.5cm, right=2.5cm, top=2.8cm, bottom=2.8cm, headheight=15pt]{geometry}
\usepackage{amsmath, amssymb}
\usepackage{booktabs}
\usepackage{array}
\usepackage{xcolor}
\usepackage{enumitem}
\usepackage{fancyhdr}
\usepackage{titlesec}
\usepackage{mdframed}
\usepackage{tabularx}
\usepackage{hyperref}

%----- Colours -----
\definecolor{sectionblue}{RGB}{10,60,110}
\definecolor{boxbg}{RGB}{255,248,240}
\definecolor{answerbg}{RGB}{240,247,255}
\definecolor{notesbg}{RGB}{255,250,230}

%----- Box environments -----
\mdfdefinestyle{qframe}{linecolor=orange!75!black, backgroundcolor=boxbg, roundcorner=4pt, linewidth=1.3pt, innertopmargin=6pt, innerbottommargin=6pt, innerleftmargin=9pt, innerrightmargin=9pt}
\mdfdefinestyle{solframe}{linecolor=sectionblue, backgroundcolor=answerbg, roundcorner=4pt, linewidth=1.3pt, innertopmargin=7pt, innerbottommargin=7pt, innerleftmargin=9pt, innerrightmargin=9pt}
\mdfdefinestyle{notesframe}{linecolor=orange!70!black, backgroundcolor=notesbg, roundcorner=4pt, linewidth=1pt, innertopmargin=5pt, innerbottommargin=5pt, innerleftmargin=8pt, innerrightmargin=8pt}
\newenvironment{qbox}{\begin{mdframed}[style=qframe]}{\end{mdframed}}
\newenvironment{solbox}{\begin{mdframed}[style=solframe]}{\end{mdframed}}
\newenvironment{notesbox}{\begin{mdframed}[style=notesframe]}{\end{mdframed}}

%----- Page style -----
\pagestyle{fancy}
\fancyhf{}
\lhead{\small\color{sectionblue} CSE 719 — Lecture 1: Distributed System Fundamentals}
\rhead{\small\color{sectionblue} Model Solutions}
\cfoot{\small\thepage}
\renewcommand{\headrulewidth}{0.5pt}

%----- Section formatting -----
\titleformat{\section}{\Large\bfseries\color{sectionblue}}{}{0em}{}[\vspace{-2pt}\color{sectionblue}\rule{\linewidth}{0.8pt}]
\titleformat{\subsection}{\large\bfseries\color{sectionblue!85!black}}{}{0em}{}

\newcommand{\qmarks}[1]{\hfill\textbf{\color{orange!70!black}[#1 marks]}}

%=====================================================================
\begin{document}
%=====================================================================
\begin{center}
  {\LARGE\bfseries\color{sectionblue} CSE 719: Distributed \& Cloud Computing}\\[0.5em]
  {\Large\bfseries Lecture 1 — Distributed System Fundamentals}\\[0.3em]
  {\large Model Answers: 2020 \& 2021 Papers}\\[0.2em]
  {\small University of Chittagong \quad|\quad 8\textsuperscript{th} Sem B.Sc.\ (Engg.) \quad|\quad Source: Lecture-01 + Tanenbaum/Sinha}
\end{center}

\vspace{0.3em}\noindent\rule{\linewidth}{1pt}\vspace{0.5em}
\tableofcontents
\newpage

%=====================================================================
\section{2020 Examination (CSE-813, 52.5 marks)}
%=====================================================================

\subsection{Q1(a) — Define distributed and cloud systems; give examples \texorpdfstring{\qmarks{2+1}}{}}
\begin{qbox}Define distributed and cloud systems. Give appropriate examples.\end{qbox}
\begin{solbox}
\textbf{Distributed System (Tanenbaum):}
``A collection of independent computers that appears to its users as a single coherent system.''
\begin{itemize}[leftmargin=1.4cm, itemsep=1pt]
  \item Each node is \textbf{autonomous, asynchronous, and failure-prone}; they communicate via an \textbf{unreliable network}.
  \item \textit{Examples:} the World Wide Web, DNS, NFS (file system), BitTorrent (P2P), Google Search.
\end{itemize}

\textbf{Cloud System:}
A model for \textbf{on-demand, network access to a shared pool of configurable computing resources} (servers, storage, applications) that are rapidly provisioned and released, billed per use.
\begin{itemize}[leftmargin=1.4cm, itemsep=1pt]
  \item \textit{Examples:} Amazon EC2/S3 (IaaS), Google App Engine (PaaS), Microsoft Office 365 (SaaS).
\end{itemize}

\textbf{Relationship:} A cloud is a large-scale distributed system delivered as a metered utility over the internet.
\end{solbox}

%-------
\subsection{Q1(b) — Tightly vs loosely coupled systems \texorpdfstring{\qmarks{3.25}}{}}
\begin{qbox}Differentiate tightly and loosely coupled systems with appropriate figures.\end{qbox}
\begin{solbox}
\begin{center}\renewcommand{\arraystretch}{1.3}
\begin{tabularx}{\linewidth}{l X X}
\toprule
\textbf{Criterion} & \textbf{Tightly Coupled} & \textbf{Loosely Coupled} \\
\midrule
Memory & \textbf{Shared} global memory & \textbf{Private} local memory per node \\
Communication & Via shared memory reads/writes & Via \textbf{message passing} over network \\
Clock & Common/synchronised clock & Independent clocks \\
Coupling & High — fast, low-latency & Low — scalable, geo-distributed \\
Failure model & One failure can affect all & Nodes fail independently \\
Example & SMP multiprocessor & Distributed system / cluster \\
\bottomrule
\end{tabularx}\end{center}

\textbf{Figures (draw in exam):}
\begin{itemize}[leftmargin=1.4cm, itemsep=2pt]
  \item \textbf{Tightly coupled:} CPU$_1$, CPU$_2$, CPU$_3$ all connected via a \underline{common bus} to \underline{one shared memory block}.
  \item \textbf{Loosely coupled:} Each node has CPU + its \underline{own memory}; nodes linked by a \underline{communication network} (LAN/WAN).
\end{itemize}
\end{solbox}

%-------
\subsection{Q1(c) — Distributed vs parallel processing \texorpdfstring{\qmarks{2.5}}{}}
\begin{qbox}In what respect are distributed computing systems better than parallel processing systems?\end{qbox}
\begin{solbox}
\begin{itemize}[leftmargin=1.4cm, itemsep=3pt]
  \item \textbf{Scalability} — can scale to geographically dispersed, virtually unlimited nodes; parallel systems are bounded by a single machine's memory bus.
  \item \textbf{Fault tolerance} — no single point of failure; a node crash need not halt the system. In shared-memory parallel systems a fault is often fatal.
  \item \textbf{Resource \& information sharing} — remote files, printers, and databases shared across sites.
  \item \textbf{Cost-effectiveness} — built from cheap commodity machines instead of expensive tightly-coupled hardware.
  \item \textbf{Geographic distribution} — services placed near users, reducing latency; parallel processing is always co-located.
\end{itemize}
\end{solbox}

%-------
\subsection{Q2(a) — Location, relocation, and migration transparency \texorpdfstring{\qmarks{3}}{}}
\begin{qbox}How are location, relocation, and migration transparencies different from each other in distributed computing? Explain with examples.\end{qbox}
\begin{solbox}
\begin{center}\renewcommand{\arraystretch}{1.35}
\begin{tabularx}{\linewidth}{l X X}
\toprule
\textbf{Type} & \textbf{What it hides} & \textbf{Example} \\
\midrule
\textbf{Location} & \textit{Where} a resource physically resides & A URL/DNS name gives no hint of the server's geographic site. \\[4pt]
\textbf{Migration} & That a resource has \textit{moved} to a new location (between client accesses) — name is unchanged & A file is silently migrated to another server; clients use the same path. \\[4pt]
\textbf{Relocation} & That a resource \textit{moved while actively in use} & A mobile user's session continues uninterrupted while the backend server changes. \\
\bottomrule
\end{tabularx}\end{center}

\textbf{Key distinction:} \textbf{Location} = you cannot see where it is. \textbf{Migration} = it moved between your accesses (name unchanged). \textbf{Relocation} = it moved \emph{during} your active connection.
\end{solbox}

%-------
\subsection{Q2(c) — Design principles for better performance \texorpdfstring{\qmarks{3}}{}}
\begin{qbox}What are the design principles considered useful for better performance in a distributed system?\end{qbox}
\begin{solbox}
\begin{enumerate}[leftmargin=1.6cm, itemsep=3pt]
  \item \textbf{Minimise communication} — batch messages, reduce round trips; exploit data locality.
  \item \textbf{Caching and replication} — keep frequently accessed data close to the consumer to reduce remote reads.
  \item \textbf{Exploit parallelism / concurrency} — overlap I/O and computation; use pipelining.
  \item \textbf{Move computation to data} — when data is large, send the program to the data node rather than the reverse.
  \item \textbf{Avoid centralised bottlenecks} — no single server or coordinator that every request must pass through; use decentralised/peer algorithms.
  \item \textbf{Load balancing} — distribute work evenly so no node is idle while others are saturated.
\end{enumerate}
\end{solbox}

%=====================================================================
\section{2021 Examination (CSE 813, 52.5 marks)}
%=====================================================================

\subsection{Q1(a) — Define a distributed system; real-time examples \texorpdfstring{\qmarks{1.75}}{}}
\begin{qbox}Define a distributed system. Mention a few real-time examples of distributed systems.\end{qbox}
\begin{solbox}
\textbf{Definition:} A collection of independent computers that appears to its users as \textbf{a single coherent system}, where entities communicate by message passing over an unreliable network.\\[4pt]
\textbf{Real-time examples:}
\begin{itemize}[leftmargin=1.4cm, itemsep=1pt]
  \item Air-traffic control systems
  \item ATM / inter-bank banking networks
  \item Online multiplayer games
  \item Sensor networks (IoT)
  \item Telecom switching networks
\end{itemize}
\end{solbox}

%-------
\subsection{Q1(b) — Advantages over standalone; disadvantages \texorpdfstring{\qmarks{2.5}}{}}
\begin{qbox}Mention the advantages of a distributed computing environment over standalone applications. What are the disadvantages of distributed systems?\end{qbox}
\begin{solbox}
\textbf{Advantages:}
\begin{itemize}[leftmargin=1.4cm, itemsep=2pt]
  \item \textbf{Resource sharing} — hardware, software, and data shared across sites.
  \item \textbf{High reliability / fault tolerance} — redundancy; no single point of failure.
  \item \textbf{Scalability} — add nodes incrementally without replacing the whole system.
  \item \textbf{Cost-effective} — commodity hardware + shared infrastructure.
  \item \textbf{Geographic distribution} — services closer to users $\Rightarrow$ lower latency.
\end{itemize}
\textbf{Disadvantages:}
\begin{itemize}[leftmargin=1.4cm, itemsep=2pt]
  \item \textbf{Complexity} — no global clock; harder to design, test, and debug.
  \item \textbf{Network dependence} — congestion/failure degrades or stops the system.
  \item \textbf{Security} — larger attack surface; data transits public/untrusted networks.
  \item \textbf{Consistency problems} — concurrent access to shared data requires synchronisation.
\end{itemize}
\end{solbox}

%-------
\subsection{Q1(c) — Scalability + challenges \texorpdfstring{\qmarks{2.5}}{}}
\begin{qbox}Define scalability. What are the challenges involved in designing scalable distributed systems?\end{qbox}
\begin{solbox}
\textbf{Scalability:} The ability of a system to \textbf{handle increasing load} (more nodes, users, or geographic spread) \textbf{without significant loss of performance}.\\
Three dimensions: \textbf{size scalability} (more nodes/users), \textbf{geographic scalability} (wider area), \textbf{administrative scalability} (more organisations).\\[4pt]
\textbf{Challenges:}
\begin{itemize}[leftmargin=1.4cm, itemsep=2pt]
  \item \textbf{Centralised components} — a single server becomes a bottleneck under high load.
  \item \textbf{Centralised data} — a single database cannot serve billions of requests.
  \item \textbf{Centralised algorithms} — any algorithm requiring global state cannot scale.
  \item \textbf{Communication latency} — grows as nodes spread across WAN distances.
  \item \textbf{Replica consistency} — keeping distributed copies in sync is harder at scale.
\end{itemize}
\textbf{Solutions:} hiding latency (async/caching), partitioning (sharding), replication.
\end{solbox}

%-------
\subsection{Q1(d) — Marshalling / Unmarshalling; ACID properties \texorpdfstring{\qmarks{2}}{}}
\begin{qbox}Define Marshalling and Unmarshalling. What are the ACID properties?\end{qbox}
\begin{solbox}
\textbf{Marshalling:} The process of \textbf{packing data items into a flat byte stream} (external representation) suitable for transmission across a network. (Also called \emph{serialisation}.)\\[3pt]
\textbf{Unmarshalling:} The inverse — \textbf{disassembling the received byte stream} back into the original data structures at the destination.\\[6pt]
\textbf{ACID Properties (transactions):}
\begin{center}\renewcommand{\arraystretch}{1.25}
\begin{tabularx}{\linewidth}{l l X}
\toprule
\textbf{Letter} & \textbf{Property} & \textbf{Meaning} \\
\midrule
A & Atomicity & All-or-nothing: a transaction either completes fully or leaves no trace. \\
C & Consistency & Takes the database from one \textbf{valid state} to another (no invariant is broken). \\
I & Isolation & Concurrent transactions do not interfere; result equals some \textbf{serial} execution. \\
D & Durability & Once committed, changes \textbf{survive crashes} (written to persistent storage). \\
\bottomrule
\end{tabularx}\end{center}
\end{solbox}

%=====================================================================
\section{Quick Summary — Exam Patterns \& Must-Know}
%=====================================================================
\begin{notesbox}
\begin{itemize}[leftmargin=1.3cm, itemsep=3pt]
  \item \textbf{Q1 of Section A = Fundamentals (2020 \& 2021).} Definition + examples open both papers — have them memorised cold.
  \item \textbf{2022 \& 2024 papers have no Lecture-1 questions.} All coverage is from 2020 and 2021.
  \item \textbf{Transparencies (2020 Q2a):} distinguish \textit{Location} vs \textit{Migration} vs \textit{Relocation} clearly — the examiner tests the differences, not just names.
  \item \textbf{Coupling (2020 Q1b):} question always says ``with appropriate figures'' — draw shared-memory vs private-memory diagrams explicitly.
  \item \textbf{Scalability (2021 Q1c):} three dimensions + the three centralisation pitfalls (component, data, algorithm) $\Rightarrow$ 5 marks in this style.
  \item \textbf{Marshalling + ACID (2021 Q1d):} two-parter for only 2 marks — keep each part to 2–3 lines.
  \item \textbf{Slide vs textbook:} DS definition/design-goals/hard-issues = Lecture-01 slides. Coupling, transparencies, ACID = Tanenbaum/Sinha textbook.
\end{itemize}
\end{notesbox}

\end{document}
